\section{Data}
\label{app:data}

This appendix documents the construction of the observable data used in the estimation in Section~\ref{sec:estimation}. All series are downloaded from the Federal Reserve Economic Data (FRED) database and converted to quarterly frequency. When FRED provides a series at monthly frequency, we take the within-quarter average.

\subsection{Raw series}
\label{app:data:fred}

Table~\ref{tab:data_fred} lists the raw FRED series used to construct the observables.

\begin{table}[h]
\centering
\caption{Raw FRED series (mnemonics) used to construct the observables for estimation and robustness exercises.}
\label{tab:data_fred}
\begin{tabular}{ll}
\hline\hline
\textbf{Observable / input} & \textbf{FRED mnemonic} \\
\hline
Real GDP (level) & GDPC1 \\
CPI price level (level) & CPIAUCSL \\
Effective federal funds rate (level, percent) & FEDFUNDS \\
Civilian non-institutional population (level) & CNP16OV \\
Civilian employment (level) & CE16OV \\
Average weekly hours (index/level) & PRS85006023 \\
GDP deflator (price level) & GDPDEF \\
Personal consumption expenditures (nominal) & PCEC \\
Private fixed investment (nominal) & FPI \\
Real government consumption expenditures and investment & GCEC1 \\
Relative price of investment goods (index) & PIRIC \\
Compensation per hour (nominal) & COMPNFB \\
\hline
\end{tabular}
\end{table}

\subsection{Transformations and scaling}
\label{app:data:transform}

The estimation uses six observables: (i) per-capita real output growth, (ii) annualized CPI inflation, (iii) the effective federal funds rate, (iv) (log) labor input, (v) per-capita real investment growth, and (vi) per-capita real government spending growth. The transformations follow the measurement conventions in \citet{smetsShocksFrictionsUS2007} and \citet{justinianoInvestmentShocksBusiness2010}.

\paragraph{Per-capita real output growth.}
Let $GDP_t$ denote real GDP (GDPC1) and $POP_t$ denote population (CNP16OV), both at quarterly frequency. Define real GDP per capita $GDP^{pc}_t \equiv GDP_t/POP_t$. The observable is quarterly percent growth,
\begin{equation}
	\Delta \ln GDP^{pc}_t \equiv 100\big(\ln GDP^{pc}_t - \ln GDP^{pc}_{t-1}\big).
\end{equation}

\paragraph{Annualized inflation.}
Let $P_t$ denote the CPI price level (CPIAUCSL). The observable inflation rate is computed as annualized log inflation,
\begin{equation}
	\pi_t \equiv 400\big(\ln P_t - \ln P_{t-1}\big).
\end{equation}

\paragraph{Federal funds rate.}
Let $FFR_t$ denote the effective federal funds rate (FEDFUNDS). This series is used in levels (percent). If downloaded at monthly frequency, it is converted to quarterly frequency by averaging within the quarter.

\paragraph{Labor input.}
Labor input combines average weekly hours (PRS85006023), employment (CE16OV), and population (CNP16OV). Let $H_t$ denote hours, $EMP_t$ employment, and $POP_t$ population. The (log) labor-input observable is
\begin{equation}
	HOURS_t \equiv 100\,\ln\Big( H_t\,\frac{EMP_t}{POP_t} \Big),
\end{equation}

and we subtract its sample mean so that $HOURS_t$ is centered at zero, consistent with the log-linear model representation.

\paragraph{Per-capita real investment growth.}
Let $INV_t$ denote nominal private fixed investment (FPI) and $DEF_t$ the GDP deflator (GDPDEF). The script constructs a real per-capita series
\begin{equation}
	INV^{pc}_t \equiv \frac{INV_t}{POP_t\,DEF_t},
\end{equation}

and then computes quarterly percent growth,
\begin{equation}
	\Delta \ln INV^{pc}_t \equiv 100\big(\ln INV^{pc}_t - \ln INV^{pc}_{t-1}\big).
\end{equation}

\paragraph{Per-capita real government spending growth.}
Let $GOV_t$ denote real government consumption expenditures and gross investment (GCEC1). The per-capita series is $GOV^{pc}_t\equiv GOV_t/POP_t$, and the observable is quarterly percent growth,
\begin{equation}
	\Delta \ln GOV^{pc}_t \equiv 100\big(\ln GOV^{pc}_t - \ln GOV^{pc}_{t-1}\big).
\end{equation}

\paragraph{Per-capita real consumption growth.}
For the SW-type robustness exercise in Section~\ref{sec:robustness}, I construct consumption growth using nominal consumption (PCEC), the GDP deflator (GDPDEF), and population (CNP16OV). Define the real per-capita consumption level
\begin{equation}
	CONS_t \equiv \frac{PCEC_t}{POP_t\,DEF_t},
\end{equation}

where $DEF_t$ denotes the GDP deflator. The observable is quarterly percent growth,
\begin{equation}
	\Delta \ln CONS_t \equiv 100\big(\ln CONS_t - \ln CONS_{t-1}\big).
\end{equation}

\paragraph{Real wage growth.}
Wage growth is based on nominal compensation per hour (COMPNFB) deflated by the GDP deflator (GDPDEF). Define the real wage level
\begin{equation}
	WAGE_t \equiv \frac{COMPNFB_t}{DEF_t},
\end{equation}

and the observable wage growth
\begin{equation}
	\Delta \ln WAGE_t \equiv 100\big(\ln WAGE_t - \ln WAGE_{t-1}\big).
\end{equation}

\paragraph{Detrended relative price of investment.}
Section~\ref{sec:robustness} uses the relative price of investment goods (PIRIC) as an additional observable to help identify the investment-specific technology (IST) shock. I first take logs and scale the series,
\begin{equation}
	RP_t \equiv 100\ln PIRIC_t,
\end{equation}

and then remove a linear trend,
\begin{equation}
	\widetilde{RP}_t \equiv \text{detrend}(RP_t).
\end{equation}

